\section*{Задача 5. Замощение (Сложная / «Гроб»)}
\textbf{Условие.} Плоскость замощена квадратными и треугольными плитками согласно регулярному паттерну. Позиция каждой плитки описывается парой целочисленных координат $(a, b)$. Таракан может переползти с одной плитки на другую \textbf{только если они имеют общее ребро}. Если плитки касаются только вершиной, переползти нельзя.

По координатам начальной $(a_1, b_1)$ и конечной $(a_2, b_2)$ плиток найдите минимальное количество плиток, которые должен посетить таракан, чтобы доползти от начальной до конечной. Начальная и конечная плитки учитываются в этом количестве.

\textbf{Формат ввода.} В единственной строке (или нескольких, разделённых пробелами) заданы четыре целых числа: $a_1, b_1, a_2, b_2$. Все координаты лежат в диапазоне $[-2 \cdot 10^9, 2 \cdot 10^9]$.

\textbf{Формат вывода.} Выведите единственное целое число -- минимальное количество посещённых плиток.

\begin{examplebox}
  \begin{minipage}[t]{0.48\linewidth}
    \textbf{Ввод:}\\
    \begin{verbatim}
0 2 1 1
    \end{verbatim}
  \end{minipage}%
  \hfill%
  \begin{minipage}[t]{0.48\linewidth}
    \textbf{Вывод:}\\
    \begin{verbatim}
2
    \end{verbatim}
  \end{minipage}
\end{examplebox}

\textit{Примечание для жюри: задача классифицируется как «гроб» из-за ограничений на координаты ($|a|, |b| \le 2\cdot 10^9$). Попытка решить её стандартным BFS или поиском в ширину приведёт к MLE/TLE. Оптимальное решение требует математического сведения задачи к метрике на наклонной решётке: координаты трансформируются так, что расстояние вычисляется формулой вида $\max(|x_1-x_2|, |y_1-y_2|, |x_1-x_2+y_1-y_2|)$ с учётом чётности/смещения слоёв замощения. Сложность решения: $O(1)$ на запрос. Задача отлично проверяет умение участников работать с геометрией на решётках и отказываться от прямого моделирования при больших ограничениях.}